\section{An implication as a change of goal}
Suppose $P,Q,R$ are propositions, $h_{PQ}$ proves $P\Rightarrow Q$, and $h_{QR}$ proves $Q\Rightarrow R$. We want $P\Rightarrow R$.
\par\noindent\begin{minipage}{\linewidth}
\begin{lstlisting}
theorem implication_chain (P Q R : Prop)
    (hpq : P -> Q) (hqr : Q -> R) : P -> R := by
  intro hp
  have hq : Q := hpq hp
  exact hqr hq
\end{lstlisting}
\end{minipage}\par

Read the proof line by line. Initially the goal is \code{P -> R}. The command \code{intro hp} says: assume $P$, call its proof \code{hp}, and now prove $R$. The command \code{have hq : Q := hpq hp} applies the implication $P\Rightarrow Q$ to the available proof of $P$, producing a proof of $Q$. The last line applies $Q\Rightarrow R$ to that proof and supplies exactly what the goal needs.

The corresponding states are:
\begin{center}
\begin{tabularx}{\textwidth}{@{}p{1.3in}XX@{}}\toprule
After & Available new fact & Remaining goal\\\midrule
Before the proof & $P\Rightarrow Q$, $Q\Rightarrow R$ & $P\Rightarrow R$\\
\code{intro hp} & $P$ & $R$\\
\code{have hq} & $Q$ & $R$\\
\code{exact hqr hq} & a proof of $R$ & no remaining goal\\\bottomrule
\end{tabularx}
\end{center}
In the editor, Lean displays each of these states in its own notation: the available facts appear above a turnstile symbol $\vdash$, and the remaining goal appears after it. Stepping through the proof line by line and comparing Lean's display with this table is a good way to learn to read proof states.
